% GENERATED FILE. Edit canonical lessons, then run npm run generate:books.
% canonical-source-hash: fnv1a64:a028c910916da628
% canonical-lessons: MR-C161-velapatrak, MR-C161-sucana, MR-C161-bhade, MR-C161-jahirat, MR-C161-pavti, MR-R161-first-pass-looking-back, MR-R161-second-pass-looking-ahead-and-reading-signs

\chapter{A Timetable, a Notice, the Rent, an Advert, a Receipt}
\label{ch:mr-a-timetable-a-notice-the-rent-an-advert-a-receipt}
\hlchaptermodality{161}

\begin{chapteropening}
\textbf{By the end of this chapter:} \emph{I can read a timetable, a notice, an advert and a receipt, and act on what they say.}

Complete the last lesson of chapter 161: I can read a timetable, a notice, an advert and a receipt, and act on what they say.
\end{chapteropening}

\section[\texorpdfstring{\mr{वेळापत्रक} (veḷāpatrak)}{वेळापत्रक (veḷāpatrak)}]{\mr{वेळापत्रक} (veḷāpatrak) — a timetable}

\label{lesson:MR-C161-velapatrak}



\begin{quote}
Before the new one: say the Marathi for from now on, then the Marathi for to plan.
\end{quote}

\subsection*{You'll want to know: \mr{वेळापत्रक}}
\textbf{\mr{वेळापत्रक}} — \emph{veḷāpatrak} — \textquotedblleft{}a timetable\textquotedblright{}. Read the board: \textbf{\mr{बस}: \mr{सकाळी} \mr{आठ} \mr{वाजता}.} The bus leaves at eight in the morning, so be there before eight.

\begin{grammarlens}[title={What you've built}]
One more word for this chapter.
\end{grammarlens}

\subsection*{Your turn}
\begin{itemize}
\raggedright
  \item \emph{Say it:} \emph{veḷāpatrak}
  \item \emph{Say it:} \emph{veḷāpatrak}, once more
  \item \emph{Say it:} say \emph{veḷāpatrak}
  \item \emph{Recall:} say the Marathi for from now on, then the Marathi for to plan, then say \emph{veḷāpatrak} again
\end{itemize}

\begin{tcolorbox}[breakable,colback=teal!4,colframe=teal!35!black,title={Before you move on}]
What does \mr{वेळापत्रक} mean? (\textquotedblleft{}A timetable\textquotedblright{}.) Say it once more.
\end{tcolorbox}
\section[\texorpdfstring{\mr{सूचना} (sūcanā)}{सूचना (sūcanā)}]{\mr{सूचना} (sūcanā) — a notice, an instruction}

\label{lesson:MR-C161-sucana}



\begin{quote}
Before the new one: say the Marathi for to plan, then the Marathi for a timetable.
\end{quote}

\subsection*{You'll want to know: \mr{सूचना}}
\textbf{\mr{सूचना}} — \emph{sūcanā} — \textquotedblleft{}a notice, an instruction\textquotedblright{}. Read the notice: \textbf{\mr{दुकान} \mr{रविवारी} \mr{बंद}.} The shop is closed on Sundays, so come on another day.

\begin{grammarlens}[title={What you've built}]
One more word for this chapter.
\end{grammarlens}

\subsection*{Your turn}
\begin{itemize}
\raggedright
  \item \emph{Say it:} \emph{sūcanā}
  \item \emph{Say it:} \emph{sūcanā}, once more
  \item \emph{Say it:} say \emph{sūcanā}
  \item \emph{Recall:} say the Marathi for to plan, then the Marathi for a timetable, then say \emph{sūcanā} again
\end{itemize}

\begin{tcolorbox}[breakable,colback=teal!4,colframe=teal!35!black,title={Before you move on}]
What does \mr{सूचना} mean? (\textquotedblleft{}A notice, an instruction\textquotedblright{}.) Say it once more.
\end{tcolorbox}
\section[\texorpdfstring{\mr{भाडे} (bhāḍe)}{भाडे (bhāḍe)}]{\mr{भाडे} (bhāḍe) — rent; a fare}

\label{lesson:MR-C161-bhade}



\begin{quote}
Before the new one: say the Marathi for a timetable, then the Marathi for a notice.
\end{quote}

\subsection*{You'll want to know: \mr{भाडे}}
\textbf{\mr{भाडे}} — \emph{bhāḍe} — \textquotedblleft{}rent; a fare\textquotedblright{}. The money for a room, or for a ride: \textbf{\mr{खोलीचे} \mr{भाडे}} — \emph{kholīce bhāḍe} — \textquotedblleft{}the rent for the room\textquotedblright{}.

\begin{grammarlens}[title={What you've built}]
One more word for this chapter.
\end{grammarlens}

\subsection*{Your turn}
\begin{itemize}
\raggedright
  \item \emph{Say it:} \emph{bhāḍe}
  \item \emph{Say it:} \emph{bhāḍe}, once more
  \item \emph{Say it:} say \emph{bhāḍe}
  \item \emph{Recall:} say the Marathi for a timetable, then the Marathi for a notice, then say \emph{bhāḍe} again
\end{itemize}

\begin{tcolorbox}[breakable,colback=teal!4,colframe=teal!35!black,title={Before you move on}]
What does \mr{भाडे} mean? (\textquotedblleft{}Rent; a fare\textquotedblright{}.) Say it once more.
\end{tcolorbox}
\section[\texorpdfstring{\mr{जाहिरात} (jāhirāt)}{जाहिरात (jāhirāt)}]{\mr{जाहिरात} (jāhirāt) — an advert}

\label{lesson:MR-C161-jahirat}



\begin{quote}
Before the new one: say the Marathi for a notice, then the Marathi for rent; a fare.
\end{quote}

\subsection*{You'll want to know: \mr{जाहिरात}}
\textbf{\mr{जाहिरात}} — \emph{jāhirāt} — \textquotedblleft{}an advert\textquotedblright{}. Read the advert: \textbf{\mr{खोली} \mr{भाड्याने} \mr{देणे} \mr{आहे}.} — \textquotedblleft{}Room to let.\textquotedblright{} A room is free to rent.

\begin{grammarlens}[title={What you've built}]
One more word for this chapter.
\end{grammarlens}

\subsection*{Your turn}
\begin{itemize}
\raggedright
  \item \emph{Say it:} \emph{jāhirāt}
  \item \emph{Say it:} \emph{jāhirāt}, once more
  \item \emph{Say it:} say \emph{jāhirāt}
  \item \emph{Recall:} say the Marathi for a notice, then the Marathi for rent; a fare, then say \emph{jāhirāt} again
\end{itemize}

\begin{tcolorbox}[breakable,colback=teal!4,colframe=teal!35!black,title={Before you move on}]
What does \mr{जाहिरात} mean? (\textquotedblleft{}An advert\textquotedblright{}.) Say it once more.
\end{tcolorbox}
\section[\texorpdfstring{\mr{पावती} (pāvtī)}{पावती (pāvtī)}]{\mr{पावती} (pāvtī) — a receipt}

\label{lesson:MR-C161-pavti}



\begin{quote}
Before the new one: say the Marathi for rent; a fare, then the Marathi for an advert.
\end{quote}

\subsection*{You'll want to know: \mr{पावती}}
\textbf{\mr{पावती}} — \emph{pāvtī} — \textquotedblleft{}a receipt\textquotedblright{}. Ask for it when you pay: \textbf{\mr{पावती} \mr{द्या}.} — \emph{pāvtī dyā} — \textquotedblleft{}A receipt, please.\textquotedblright{}

\begin{grammarlens}[title={What you've built}]
One more word for this chapter.
\end{grammarlens}

\subsection*{Your turn}
\begin{itemize}
\raggedright
  \item \emph{Say it:} \emph{pāvtī}
  \item \emph{Say it:} \emph{pāvtī}, once more
  \item \emph{Say it:} say \emph{pāvtī}
  \item \emph{Recall:} say the Marathi for rent; a fare, then the Marathi for an advert, then say \emph{pāvtī} again
\end{itemize}

\begin{tcolorbox}[breakable,colback=teal!4,colframe=teal!35!black,title={Before you move on}]
What does \mr{पावती} mean? (\textquotedblleft{}A receipt\textquotedblright{}.) Say it once more.
\end{tcolorbox}
\section[(dialogue)]{Looking back — a first pass}

\label{lesson:MR-R161-first-pass-looking-back}



\begin{quote}
Say the words for an advert and a receipt.
\end{quote}

\subsection*{Your turn}
\begin{itemize}
\raggedright
  \item \emph{Say it:} recently, formerly, already, the past, to happen
\end{itemize}

\begin{tcolorbox}[breakable,colback=teal!4,colframe=teal!35!black,title={Before you move on}]
Recently? (\textbf{\mr{अलीकडे}}, \emph{alīkaḍe}.) Formerly? (\textbf{\mr{पूर्वी}}, \emph{pūrvī}.) Already? (\textbf{\mr{आधीच}}, \emph{ādhīc}.) The past? (\textbf{\mr{भूतकाळ}}, \emph{bhūtkāḷ}.) To happen? (\textbf{\mr{घडणे}}, \emph{ghaḍṇe}.)
\end{tcolorbox}
\section[(dialogue)]{Looking ahead and reading signs — a second pass}

\label{lesson:MR-R161-second-pass-looking-ahead-and-reading-signs}



\begin{quote}
Say the words for the future and soon.
\end{quote}

\subsection*{Your turn}
Say these aloud:

\begin{itemize}
\raggedright
  \item the future, soon, next, from now on, to plan
  \item a timetable, a notice, rent; a fare, an advert, a receipt
\end{itemize}

\begin{tcolorbox}[breakable,colback=teal!4,colframe=teal!35!black,title={Before you move on}]
The future? (\textbf{\mr{भविष्य}}, \emph{bhaviṣya}.) Soon? (\textbf{\mr{लवकरच}}, \emph{lavkarac}.) Next? (\textbf{\mr{पुढचा} / \mr{पुढची} / \mr{पुढचे}}, \emph{puḍhcā / puḍhcī / puḍhce}.) From now on? (\textbf{\mr{यानंतर}}, \emph{yānantar}.) To plan? (\textbf{\mr{आखणे}}, \emph{ākhṇe}.) A timetable? (\textbf{\mr{वेळापत्रक}}, \emph{veḷāpatrak}.) A notice? (\textbf{\mr{सूचना}}, \emph{sūcanā}.) Rent; a fare? (\textbf{\mr{भाडे}}, \emph{bhāḍe}.) An advert? (\textbf{\mr{जाहिरात}}, \emph{jāhirāt}.) A receipt? (\textbf{\mr{पावती}}, \emph{pāvtī}.)
\end{tcolorbox}
